% Generated draft: requires mathematical and source review.
\documentclass[11pt]{article}
\usepackage[margin=25mm]{geometry}
\usepackage{fontspec}
\setmainfont{Latin Modern Roman}
\usepackage{amsmath,amssymb,mathtools,unicode-math}
\setmathfont{Latin Modern Math}
\usepackage{enumitem,needspace,xurl,hyperref}
\setlength{\parindent}{0pt}
\setlength{\parskip}{5pt}
\setlength{\emergencystretch}{3em}
\newcommand{\N}{\mathbb{N}}
\newcommand{\Z}{\mathbb{Z}}
\newcommand{\Q}{\mathbb{Q}}
\newcommand{\R}{\mathbb{R}}
\newcommand{\C}{\mathbb{C}}
\newcommand{\F}{\mathbb{F}}
\newcommand{\A}{\mathbb{A}}
\newcommand{\PP}{\mathbb{P}}
\newcommand{\E}{\mathbb{E}}
\newcommand{\eps}{\varepsilon}
\newcommand{\dd}{\,\mathrm{d}}
\begin{document}
\section*{UnsolvedMath ID: 30000962}
\textbf{Source problem code:} catalogue entry
\subsection{Definitions and mathematical statement}
The principal statement is Conjecture 11. Let $\mathfrak g$ be a complex reductive Lie algebra, $K$ a knot, $P$ the intended weight lattice, and $W$ its Weyl group. Let $J_{\mathfrak g,K}:P\to R_q$ be the source's Witten--Reshetikhin--Turaev function, with its intended coefficient ring $R_q$. Let $I^W_{\mathfrak g,K}$ be its Weyl-invariant recursive ideal acting on weight-lattice functions by the prescribed recurrence operators. These coefficient, action and normalization data are not supplied.

The conjecture asks whether, for every such $\mathfrak g$ and $K$, there exists a finite subset $S\subset P$ such that
\[
\left.
\begin{gathered}
F:P\to R_q\text{ is Weyl-equivariant},\qquad
D F=0\quad(D\in I^W_{\mathfrak g,K}),\\
F(\lambda)=J_{\mathfrak g,K}(\lambda)\quad(\lambda\in S)
\end{gathered}
\right\}
\quad\Longrightarrow\quad F=J_{\mathfrak g,K}.
\]
The finite set may depend on $\mathfrak g,K$ but must work for all competing $F$. No uniform bound on its size or algorithm to find it is asserted. Weyl-equivariance is retained as a source-dependent condition: ordinary invariance, a shifted Weyl action and a sign-twisted action must not be silently identified.

The longer original field also includes a character-variety assertion. The verification note explicitly says that this is the separate Conjecture 12, represented by OWR-1967-010, and that the present record now contains Conjecture 11 only. To preserve that source context without merging the two, the additional original assertion is recorded here as follows. Let $G$ be the intended complex group associated with $\mathfrak g$, $M_K$ the knot complement and $r$ boundary restriction. Writing $X_G(M)=\operatorname{Hom}(\pi_1M,G)//G$ in the standard character-variety convention, its assertion is
\[
V\bigl(\varepsilon(I^W_{\mathfrak g,K})\bigr)
=\overline{r(X_G(M_K))}^{\,\mathrm{Zar}}
\subset X_G(\partial M_K),\qquad
\sqrt{\varepsilon(I^W_{\mathfrak g,K})}=A_G(K),
\]
where $\varepsilon$ specializes $q$ to $1$ and $A_G(K)$ is the vanishing ideal of that image closure. The equivalence uses the appropriate commutative coordinate-ring and algebraic-closure conventions. The choice of $G$, the specialization and ideal identifications are not provided. This is attributed context for the separately indexed conjecture, not a second primary claim assigned to this record.

The supplied triage checked 2026-08-21 labels Conjecture 11 open and says no general proof or disproof was located. This is a dated source assessment.
\subsection{Short English statement}
Ask whether finitely many values and the Weyl-invariant recurrence ideal uniquely specify the quantum knot function among Weyl-equivariant functions. The older character-variety clause is retained as context explicitly assigned by source verification to the separate Conjecture 12.
\subsection{Sources}
\begin{itemize}
\item Local UnsolvedMath registry, permanent ID 30000962, source code catalogue entry; supplied complete record and its corrected-verified note separating Conjectures 11 and 12. The latter is cross-referenced there as OWR-1967-010; no permanent numeric ID for it is supplied here.
\item Invariants in Low-Dimensional Topology (2009), Oberwolfach Reports, DOI 10.4171/owr/2008/22, pp. 1157--1232: \url{https://doi.org/10.4171/owr/2008/22}. Proposed-year metadata gives 2008.
\end{itemize}
\end{document}
